\section{Gli organismi viventi}

Gli organismi viventi sono in grado di:
\begin{itemize}
  \item riprodursi e trasmettere i propri caratteri (eredità);
  \item compiere un ciclo vitale;
  \item rinnovare continuamente la propria struttura;
  \item reagire agli stimoli;
  \item muoversi;
  \item evolversi.
\end{itemize}

% ---------------------------------------------------------------------------
\subsection{Riproduzione}

Insieme dei meccanismi mediante i quali gli esseri viventi provvedono alla conservazione della propria specie, generando nuovi individui simili a sé che subentreranno al genitore, o ai genitori, nella popolazione.

\rubrica{Riproduzione asessuata}
Genera individui che mantengono invariato il patrimonio genetico del genitore.
\begin{itemize}
  \item \textbf{moltiplicazione cellulare} (divisione mitotica): una cellula madre si divide in due cellule figlie identiche; negli organismi unicellulari è l'unica forma di riproduzione e incrementa una popolazione geneticamente uguale, nei pluricellulari è il meccanismo dell'accrescimento del singolo individuo;
  \item \textbf{scissione}: un individuo generante, dividendo il proprio corpo, genera individui figli identici tra loro e al generante;
  \item \textbf{frammentazione}: una parte dell'organismo si stacca e forma un nuovo individuo completo e identico al generante (per rigenerazione);
  \item \textbf{gemmazione}: formazione di gemme laterali che si separano definitivamente mediante una strozzatura, dando origine a nuovi individui uguali al generante;
  \item \textbf{poliembrionia}: divisione in due o più parti dello zigote o dell'embrione nei primissimi stadi di sviluppo $\rightarrow$ gemelli monozigoti;
  \item \textbf{sporulazione}: cellule specializzate (sporocisti), in condizioni ambientali sfavorevoli, producono per divisione mitotica speciali cellule riproduttive (\textit{mitospore}) capaci di generare un nuovo individuo quando il contesto ridiventa favorevole; le mitospore hanno una spessa parete protettiva che permette di resistere alle condizioni avverse.
\end{itemize}

\rubrica{Riproduzione sessuata}
Il nuovo individuo si origina dalla fusione dei gameti; la divisione che produce i gameti è la \textbf{meiosi}. Aumenta la variabilità genetica all'interno di una specie.
\begin{itemize}
  \item \textbf{anfigonia}: fecondazione di due gameti per formare una nuova cellula, lo zigote; la fusione di due gameti aploidi genera una cellula diploide (lo zigote), con metà del materiale genetico dal gamete maschile e metà da quello femminile;
  \item \textbf{partenogenesi}: riproduzione tramite sviluppo di gameti femminili in assenza di fecondazione; avviene generalmente per auto-attivazione dell'uovo non fecondato che, per restituzione anafasica, ricostituisce un genoma diploide omozigote in tutti i loci (due assetti aploidi identici). Generalmente è un evento accidentale, ma in alcune specie è una forma alternativa all'anfigonia e genera prole: es.\ \textit{Bacillus rossius} (insetto stecco) e alcuni imenotteri sociali (api, vespe).
\end{itemize}

% ---------------------------------------------------------------------------
\subsection{Ciclo vitale}

$\text{nascita} \rightarrow \text{crescita} \rightarrow \text{sviluppo} \rightarrow \text{morte}$; durata variabile a seconda della specie.
\begin{itemize}
  \item \textbf{crescita}: aumento della dimensione cellulare e/o del numero di cellule;
  \item \textbf{sviluppo}: modificazioni strutturali e funzionali, fino a un progressivo rallentamento e deterioramento delle funzionalità che porta alla morte.
\end{itemize}

% ---------------------------------------------------------------------------
\subsection{Metabolismo}

Il complesso delle trasformazioni di natura chimica che avvengono negli organismi viventi.
\begin{itemize}
  \item \textbf{capacità di nutrirsi}: organismi autotrofi, organismi eterotrofi;
  \item \textbf{capacità di trasformare materia ed energia} (insieme delle reazioni chimiche e fisiche che avvengono in un organismo o in una sua parte):
    \begin{itemize}
      \item \textit{anabolismo}: produce molecole complesse a partire da molecole più semplici, con consumo di energia (reazioni endoergoniche); es.\ biosintesi di una proteina;
      \item \textit{catabolismo}: degrada molecole complesse in molecole più semplici (reazioni esoergoniche), liberando energia; es.\ respirazione cellulare o ciclo di Krebs.
    \end{itemize}
\end{itemize}

% ---------------------------------------------------------------------------
\subsection{Reazione agli stimoli}

Capacità di percepire e reagire a stimoli esterni: cambiamenti di temperatura e di pressione, presenza/assenza di luce, cambiamenti chimici. Esempi: formazione di spore, letargo, mimetismo, reazione ai rumori o a stimoli olfattivi, fototropismo.

% ---------------------------------------------------------------------------
\subsection{Movimento}

\begin{itemize}
  \item \textbf{organismi unicellulari}: movimento ameboide (caratteristico delle \textit{Amoebe}, che si muovono emettendo prolungamenti detti pseudopodi); oppure mediante ciglia o flagelli;
  \item \textbf{organismi superiori}:
    \begin{itemize}
      \item \textit{piante}: allungamenti o rotazioni; movimenti lenti e limitati, in funzione della presenza di acqua (radici) o del sole (fronde, fiori);
      \item \textit{animali}: organi preposti al movimento, che può essere rapidissimo e coprire notevoli distanze (pinne, ali, zampe, arti).
    \end{itemize}
\end{itemize}

% ---------------------------------------------------------------------------
\subsection{Ambienti biologici}

\begin{itemize}
  \item \textbf{ambiente acqueo}: marino, d'acqua dolce;
  \item \textbf{ambiente terrestre}: epigeo, ipogeo.
\end{itemize}

\rubrica{Ambiente marino}
\begin{itemize}
  \item \textbf{dominio bentonico} (o di fondo): ambiente dove vivono gli organismi legati più o meno direttamente ai fondali:
    \begin{itemize}
      \item \textit{ambiente litoraneo}: acque e fondali che circondano le coste (la cosiddetta \textit{platea continentale}), fino a una profondità massima di 200 m; molto eterogeneo (in superficie risente dei moti ondosi ed è ben illuminato, in profondità è quasi buio e il moto ondoso quasi assente); i livelli meno profondi sono i più ricchi di vita animale e vegetale;
      \item \textit{ambiente profondo}: oltre la platea continentale il mare sprofonda più rapidamente (a gradino);
    \end{itemize}
  \item \textbf{dominio pelagico}: acque libere, distanti dalle coste e dal fondo, che si estendono dalla superficie fino agli abissi delle fosse oceaniche; vi vivono gli organismi che conducono una vita non vincolata in maniera esclusiva al fondale.
\end{itemize}
Nell'ambiente marino la variabilità maggiore è data da salinità e ossigenazione.

\rubrica{Organismi dell'ambiente marino}
\begin{itemize}
  \item \textbf{plancton}: organismi animali e vegetali che galleggiano, trasportati quasi tutti passivamente da moti ondosi e correnti (alghe unicellulari, protozoi, larve, piccoli crostacei, meduse, alghe pluricellulari);
  \item \textbf{necton}: organismi che nuotano attivamente, provvisti di mezzi per nuotare, emergere e immergersi (pesci, cefalopodi, tartarughe, cetacei);
  \item \textbf{benthos}: organismi che vivono in stretto contatto con il fondo o fissati a un substrato solido (alghe pluricellulari, animali che camminano o strisciano come i crostacei, animali sessili ancorati a un substrato e incapaci di muoversi come i molluschi bivalvi).
\end{itemize}

\rubrica{Ambienti d'acqua dolce}
\begin{itemize}
  \item \textbf{laghi}: ambiente in parte sovrapponibile a quello marino; i laghi di montagna, più freddi, hanno flora e fauna limitata rispetto a quelli di pianura;
  \item \textbf{fiumi, torrenti e sorgenti}: flora e fauna diverse a seconda del corso (lento o impetuoso) e delle acque (calde o fredde).
\end{itemize}

\rubrica{Ambiente terrestre}
Varia in base a latitudine, longitudine, sottosuolo, vegetazione, clima e stagioni; variabilità maggiore data da temperatura e umidità.
\begin{itemize}
  \item \textbf{epigeo}: organismi che vivono sulla superficie della terra;
  \item \textbf{ipogeo}: organismi che vivono dentro la terra (grotte, anfratti, gallerie scavate nel sottosuolo);
  \item \textbf{entozoico}: endoparassiti che vivono all'interno di un altro animale ospite.
\end{itemize}

% ---------------------------------------------------------------------------
\subsection{Rapporti intraspecifici}

Rapporti che intercorrono tra individui della stessa specie.
\begin{itemize}
  \item \textbf{colonie}: individui della stessa specie che, dopo essersi riprodotti per via asessuata (gemmazione), non si separano e rimangono fisicamente uniti:
    \begin{itemize}
      \item \textit{colonie omeomorfe}: tutti gli individui sono uguali (spugne, madrepore, coralli);
      \item \textit{colonie eteromorfe}: gruppi di individui si specializzano in una particolare funzione fino a una spiccata differenziazione morfo-funzionale (``superorganismi'');
    \end{itemize}
  \item \textbf{società animali}: rapporti tra animali della stessa specie che, pur fisicamente separati, conducono vita comune (temporanee, durature, individualiste, collettiviste).
\end{itemize}

% ---------------------------------------------------------------------------
\subsection{Rapporti interspecifici}

Rapporti tra individui di specie diversa che si instaurano all'interno di uno stesso ambiente; riguardano la nutrizione e l'occupazione degli spazi.

\rubrica{Simbiosi}
Stretta relazione fra individui di specie diverse che convivono per trarre un beneficio reciproco.
\begin{itemize}
  \item \textbf{mutualismo}: nessuno degli individui è in grado di vivere isolato (licheni; paguro e attinia; \textit{Rhizobium});
  \item \textbf{commensalismo}: solo una delle due specie trae vantaggio dall'associazione, senza però danneggiare l'altra (pesce pagliaccio; acari foretici; storno);
  \item \textbf{inquilinismo}: i rapporti fra le due specie si riducono all'occupazione di uno spazio comune.
\end{itemize}

\rubrica{Amensalismo}
Una specie impedisce o diminuisce il successo di un'altra, senza però trarne vantaggio. Si ha quando un organismo produce una sostanza chimica, come parte del suo normale metabolismo, con effetto negativo su altri organismi (es.\ \textit{Penicillium}, che secerne penicillina, un potente battericida).

\rubrica{Parassitismo}
Una delle due specie conviventi, il parassita, trae vantaggio dall'associazione a scapito dell'altra, l'ospite, creandole un danno biologico.
\begin{itemize}
  \item il parassita dipende dall'ospite, a cui è più o meno intimamente legato da una relazione anatomica e fisiologica;
  \item il parassita ha una struttura anatomica più semplice di quella dell'ospite;
  \item l'ospite è più grande del parassita;
  \item il ciclo vitale del parassita è più breve di quello dell'ospite e si conclude prima della sua morte;
  \item il parassita ha rapporti con un solo ospite, ma non viceversa.
\end{itemize}
Tipi:
\begin{itemize}
  \item \textbf{facoltativo}: non si hanno modificazioni morfo-funzionali;
  \item \textbf{obbligato}: la vita del parassita è subordinata a quella dell'ospite;
  \item \textbf{ectoparassiti}: adottano elaborati meccanismi e strategie per attaccare un ospite (es.\ le sanguisughe individuano l'ospite tramite sensori di movimento e ne accertano l'identità attraverso la temperatura della pelle e indicazioni chimiche; hanno un adattamento morfofunzionale per aderire all'ospite e succhiarne il nutrimento);
  \item \textbf{endoparassiti}: attaccano l'ospite in modo passivo (es.\ gli endoparassiti dell'intestino umano infestano l'uomo tramite l'ingestione di acqua sporca o cibi infestati; hanno subìto grosse modificazioni morfofunzionali).
\end{itemize}
Esempi: \textit{Taenia solium}, \textit{Plasmodium malariae}.
